% from ega3-6-fr.tex
\oldpage[III]{157}
On peut considérer que les isomorphismes (6.6.7, $t$) définissent des systèmes
inductifs de foncteurs spectraux sur l'ensemble filtrant des couples
$(\mathfrak U^{(1)},\mathfrak U^{(2)})$ de recouvrements ouverts affines finis
de $X^{(1)}$ et $X^{(2)}$. Nous désignerons par
\[
  {}^{(t)}\mathscr E
  (X^{(1)},X^{(2)};
   \mathscr P_\bullet^{(1)},\mathscr P_\bullet^{(2)})
  \quad\text{ou}\quad
  {}^{(t)}\mathscr E^S
  (X^{(1)},X^{(2)};
   \mathscr P_\bullet^{(1)},\mathscr P_\bullet^{(2)})
\]
la limite inductive de ce système, par
$\mathbf{Tor}_\bullet^S
(X^{(1)},X^{(2)};\mathscr P_\bullet^{(1)},\mathscr P_\bullet^{(2)})$
l'aboutissement de ce foncteur spectral, que nous appellerons l'\emph{hypertor
global} des deux complexes $\mathscr P_\bullet^{(1)}$ et
$\mathscr P_\bullet^{(2)}$; si $\mathscr P_\bullet^{(1)}$ et
$\mathscr P_\bullet^{(2)}$ sont réduits à leurs termes de degré 0,
$\mathscr F^{(1)}$ et $\mathscr F^{(2)}$, nous écrirons
\[
  \mathbf{Tor}_\bullet^S
  (X^{(1)},X^{(2)};\mathscr F^{(1)},\mathscr F^{(2)}),
\]
et conformément aux conventions générales,
\[
  \mathbf{Tor}_q^S
  (X^{(1)},X^{(2)};\mathscr P_\bullet^{(1)},\mathscr P_\bullet^{(2)})
\]
sera donc le bicomplexe des
\[
  \mathbf{Tor}_q^S
  (X^{(1)},X^{(2)};\mathscr P_h^{(1)},\mathscr P_k^{(2)}).
\]
\end{env}

\begin{env}[6.6.8]
Les hypothèses étant celles de (6.6.1), considérons maintenant deux
$S$-morphismes $f_i:X^{(i)}\to Y^{(i)}$, où
$Y^{(i)}=\operatorname{Spec}(B_i)$ est un $S$-schéma \emph{affine}, $B_i$
étant donc une $A$-algèbre ($i=1,2$); cela définit donc un $A$-homomorphisme
$B_i\to\Gamma(X^{(i)},\mathscr O_{X^{(i)}})$ (I, 2.2.4), et par suite chacun
des $L_{\bullet\bullet}^{(i)}$ définis dans (6.6.2) est un bicomplexe de
$B_i$-modules; on en conclut que
$L_{\bullet\bullet}^{(1)}\otimes_A L_{\bullet\bullet}^{(2)}$ est un
quadricomplexe de $(B_1\otimes_A B_2)$-modules, et ses six foncteurs spectraux
d'hyperhomologie peuvent donc être considérés comme prenant leurs valeurs
dans la catégorie des suites spectrales de $(B_1\otimes_A B_2)$-modules. Si
l'on pose
$Y=Y^{(1)}\times_S Y^{(2)}=\operatorname{Spec}(B_1\otimes_A B_2)$, on peut
considérer les $\mathscr O_Y$-Modules quasi-cohérents associés à ces modules
(I, 1.3.4); nous noterons
${}^{(t)}\mathscr E
(f_1,f_2;\mathscr P_\bullet^{(1)},\mathscr P_\bullet^{(2)})$
(pour $t=a$, $b$, $a'$, $b'$, $c$ ou $d$) les six suites spectrales
\[
  \bigl({}^{(t)}\mathscr E
  (X^{(1)},X^{(2)};
   \mathscr P_\bullet^{(1)},\mathscr P_\bullet^{(2)})\bigr)^{\sim}
\]
de $\mathscr O_Y$-Modules, et
$\mathfrak{Tor}_\bullet^S
(f_1,f_2;\mathscr P_\bullet^{(1)},\mathscr P_\bullet^{(2)})$ leur
aboutissement commun
\[
  \bigl(\mathbf{Tor}_\bullet^S
  (X^{(1)},X^{(2)};
   \mathscr P_\bullet^{(1)},\mathscr P_\bullet^{(2)})\bigr)^{\sim}.
\]
On le notera
$\mathfrak{Tor}_\bullet^S(f_1,f_2;\mathscr F^{(1)},\mathscr F^{(2)})$
lorsque $\mathscr P_\bullet^{(i)}$ est réduit à son terme de degré 0,
$\mathscr F^{(i)}$ ($i=1,2$).
\end{env}

\subsection{Foncteurs hypertor globaux de complexes de Modules quasi-cohérents
et suites spectrales de Künneth : cas général.}

\begin{env}[6.7.1]
Nous allons maintenant généraliser les définitions de (6.6.8) au cas où $S$
est un préschéma quelconque, $X^{(i)}$, $Y^{(i)}$ des $S$-préschémas et
$f_i:X^{(i)}\to Y^{(i)}$ des morphismes \emph{séparés et quasi-compacts}. Il
s'agit alors, pour tout couple de complexes $\mathscr P_\bullet^{(i)}$ de
$\mathscr O_{X^{(i)}}$-Modules quasi-cohérents, limités inférieurement
($i=1,2$), de définir pour tout $n$ un $\mathscr O_Y$-Module quasi-cohérent
$\mathfrak{Tor}_n^S
(f_1,f_2;\mathscr P_\bullet^{(1)},\mathscr P_\bullet^{(2)})$ ainsi que 6
foncteurs spectraux, se réduisant aux définitions de (6.6.8) lorsque $S$,
$Y^{(1)}$ et $Y^{(2)}$ sont affines (on pose
$Y=Y^{(1)}\times_S Y^{(2)}$).

Supposons d'abord $S=\operatorname{Spec}(A)$ \emph{affine}, mais $Y^{(1)}$
et $Y^{(2)}$ quelconques; soit $W^{(i)}$ un ouvert affine de $Y^{(i)}$;
$f_i^{-1}(W^{(i)})$ est alors un $S$-schéma quasi-compact,
$W=W^{(1)}\times_S W^{(2)}$ un ouvert affine de $Y$; soit
$f'_i:f_i^{-1}(W^{(i)})\to W^{(i)}$ la restriction de $f_i$, et
$\mathscr P'_\bullet{}^{(i)}$ la restriction
$\mathscr P_\bullet^{(i)}|f_i^{-1}(W^{(i)})$ ($i=1,2$). On a alors d'après
(6.6.8) les suites spectrales
${}^{(t)}\mathscr E
(f'_1,f'_2;\mathscr P'_\bullet{}^{(1)},\mathscr P'_\bullet{}^{(2)})$
de $(\mathscr O_Y|W)$-Modules quasi-cohérents, et il s'agit de vérifier
qu'elles satisfont aux conditions de recollement ($0_{\mathrm I}$, 3.3.1).
On est aussitôt ramené au cas où $Y^{(i)}=\operatorname{Spec}(B_i)$ est affine
et où $W^{(i)}=D(g_i)$, où $g_i\in B_i$, de sorte que
$W=D(g_1\otimes g_2)$
